\section{Conclusion}
\label{sec:conclusion}
% TARGET LENGTH: 0.75--1 page

We presented Kaan, an offline phone system for four-class acoustic screening of stored-rice pests, together with a leakage-aware multi-approach evaluation.
Under file-level, byte-deduped splits of IRRI pest audio and ambient clean windows, several classical models and a strong mel-CNN exceed the cited 84.51\% reference line, with gradient boosting and the deep mel-CNN statistically compatible on a paired seed-42 fold.
Training recipe ablations show that deep models can fail without careful optimization, while error analysis highlights the weevil--borer confusion as the main remaining difficulty.
The immediate scientific next step is a phone-on-bag external test set collected in Indian storage settings; until then, Kaan should be understood as an open screening tool with workshop-grade evidence---not a field-validated diagnostic.
